\documentclass[12pt]{article}
\usepackage[utf8]{inputenc}
\usepackage[T1]{fontenc}
\usepackage{hyperref}
\title{Timucua Language Resources for K-12 Curriculum Development}
\author{Compiled by Claude Sonnet 5 (claude-sonnet-5) for Andr\'as Kornai}
\date{}
\begin{document}
\maketitle

\section{Introduction}
Timucua (ISO 639-3 tjm, Glottocode timu1245) is an extinct language isolate of the southeastern United States, with no native speakers. Most of the documentation comes from the Mocama and Potano dialects. It is known almost entirely from Spanish colonial missionary texts, chiefly those of Francisco Pareja. Any K-12 curriculum would therefore be a revitalization or heritage curriculum built on reconstruction from these texts. I found no ready-made school curriculum in the search results.

\section{Primary Educational Resources}
I found no primers, textbooks or classroom materials for Timucua in the search results. Pareja's 1612 catechism was the first book published in an indigenous language of what is now the United States. The parallel Spanish--Timucua religious texts could serve as source material for teachers, but only after adaptation.

\section{Linguistic and Reference Materials}
\begin{itemize}
\item Overview of the language: \url{https://en.wikipedia.org/wiki/Timucua_language}
\item Broadwell's online Timucua lexicon project (Webonary): \url{https://www.webonary.work/timucua/?p=9444} and \url{https://www.webonary.work/timucua/?p=146853}
\item Article in the North Florida Journal of Anthropology: \url{https://journals.flvc.org/NFJA/article/download/123622/128600/213276}
\item Broadwell, paper on intention and authorship in Timucua texts: \url{https://people.clas.ufl.edu/broadwell/files/Intention-and-authorship-revised-to-meet-editors-requests.pdf}
\item Academia.edu listing of Timucua scholarship: \url{https://www.academia.edu/Documents/in/Timucua_language}
\end{itemize}
The search summary reported that George Aaron Broadwell published \emph{Timucua: A Text-Based Reference Grammar} (University of Nebraska Press, 2024). I did not open a page for the book itself, so it has no URL here. The summary also reported that no colonial-era Timucua dictionary is known, and that Pareja's \emph{Arte} (1614) is the clearest lexical source.

\section{Religious and Community Resources}
The surviving corpus is mostly Catholic devotional material by Pareja: catechisms, a confessional and a doctrina. The Library of Congress blog discusses writing Timucua in seventeenth-century Florida: \url{https://blogs.loc.gov/maps/2021/07/chi-mobi-writing-timucua-in-seventeenth-century-florida/}. I did not find a Timucua community organization in the results.

\section{Digital and Audio Resources}
The Webonary lexicon is the only digital resource I found. I found no audio recordings. This is expected, because the language predates recording and has no speakers.

\section{Curriculum Development Context}
Little exists for classroom use. A realistic heritage curriculum could cover the history of the language, the colonial documents, and simple vocabulary drawn from the Webonary lexicon. The reconstructed pronunciation and grammar are uncertain, and a curriculum should say so. Consult Broadwell's grammar and current scholars before building lessons.
\end{document}